\chapter{引言}\label{chap:introduction}
\SourcePage{24}

镜像对称在代数几何中作出了一些令人惊讶的预测，所涉及的对象从五次三维簇上的有理曲线数目，到某些模空间的结构。这些都是极有价值的研究问题，而且如前言所述，已经催生了许多非常有趣的数学。然而，若要弄清这些预测从何而来，代数几何学家就必须深入物理学的语言；这种语言既陌生，有时又因欠缺严格性而令人沮丧。因此，在开始考察镜像对称的代数几何之前，我们先讨论促成镜像对称的动机，以及镜像对称在物理学中意味着什么。由于会涉及许多读者可能不熟悉的物理学术语，我们的介绍难免不够完整。尽管如此，我们仍希望传达这一非凡现象背后的一些直观想法。

接着，我们将讨论三点函数——它们对于镜像对称的枚举预测至关重要——以及 Calabi--Yau 流形出现在这一理论中的物理原因。最后，在本章结尾，我们会回到读者更熟悉的数学世界，预览本书其余章节将要探索的代数几何。

\section{镜像对称的物理学}\label{sec:physics}

本节旨在让读者体会镜像对称为什么会出现，以及它应当蕴含什么。从物理学的观点看，镜像对称自然地产生于超对称弦理论的标准构造；因此，我们会先对弦和超对称作一些初步说明。请读者放心：我们不预设任何物理学背景！我们的目的，是传达这些物理理论的风貌，并借此增强读者对相应数学的直观认识。读者不必详尽掌握其中的物理学，虽然后面的章节会引用本章介绍的物理内容。关于弦理论的一般参考资料，可参阅 \SourceCite{GSW,Polchinski2}。\index{超对称弦理论|hyperpage}

在弦理论中，物理过程由弦在时空中的传播来描述。传播中的弦扫出一个曲面，称为弦的\emph{世界面} $\Sigma$。经典场可以描述为世界面上的函数、丛的截面等；将它们量子化，便得到一个二维量子场论。这个理论具有一个广义的态 Hilbert 空间，也有一组可观测量；后者成为态空间上的自伴算子。理论的另外两个关键组成部分，是把 Lagrange 量在世界面 $\Sigma$ 上积分得到的作用量 $S$，以及\emph{关联函数}
\begin{equation}\label{eq:correlation}
 \bigl\langle\phi(x_1),\ldots,\phi(x_n)\bigr\rangle
 =\int[D\phi]\,\phi(x_1)\cdots\phi(x_n)e^{iS(\phi)}.
\end{equation}
\SourcePage{25}
这里 $\phi$ 是一个可观测量。这个 Feynman 积分遍及所有可能的世界面；到目前为止，它在数学上尚无定义。\hyperref[app:physical-theories]{附录 B} 会更详细地讨论这类理论。现在，我们先关注它们的一般特征。

弦理论有一个优点：它消除了粒子分裂为两个粒子时出现的某些问题。下图表明，如果把粒子表示为点，就会出现一个奇点；用弦来表示时，得到的则是一个光滑的带边二维流形：
\begin{figure}[H]
 \centering\includegraphics[width=0.79219\textwidth]{1-1-string-splitting.png}
\end{figure}

不过，弦理论仍有一些不理想的特征，其中包括许多需要重整化处理的无穷大。近来一个非凡的发现是，超对称能够消除其中许多困难。超对称把玻色子——自旋为整数、波函数对称的粒子——变为费米子——自旋为半整数、波函数反对称的粒子——反之亦然。虽然迄今尚无实验证实超对称，但由于超对称理论具有良好的性质，它们已在理论物理学中变得十分重要。\index{超对称|hyperpage}

我们还需要另一个要素：世界面 $\Sigma$ 带有共形结构，我们的超对称弦理论必须在共形等价下保持等价。因此，该理论是一种\emph{超共形场论}，简称 SCFT。这类理论的对称群的 Lie 代数称为超共形代数。它以共形代数——世界面共形变换群的 Lie 代数——为子代数，同时还包含超对称变换。\index{超共形场论|hyperpage}

\IndexLocation{indexloc:1306:25}\IndexLocation{indexloc:1307:25}\IndexLocation{indexloc:1308:25}超弦理论有四种基本类型：I 型、IIA 型、IIB 型，以及我们最感兴趣的杂化弦理论。杂化弦理论是一种 $N=2$ 的 SCFT，因为费米子解耦为左行解和右行解；这意味着实际上有四种超对称，其中两种左行、两种右行。因此，更准确地说，杂化弦理论是一种 $(2,2)$ 理论：其左行部分与右行部分各有两种独立的超对称。

$N=2$ 超共形代数包含通常的超共形代数的两个副本，因此有两个 $u(1)$ 子代数。一个位于理论的右行部分，对两种超对称施加无穷小旋转；另一个位于左行部分，作用方式相似。给两种超对称变换排定顺序，便可为每个 $u(1)$ 选取一个生成元，但这种选取并非典范的。如果颠倒两种超对称的顺序，$u(1)$ 生成元的符号就会改变。这两个子代数各自的生成元记为
\SourcePage{26}
$(Q,\overline Q)$，它们除符号外唯一确定。我们把 $(Q,\overline Q)$ 看作作用于态 Hilbert 空间的算子；在 $u(1)\times u(1)$ 的作用下，该空间分解为特征子空间。正如即将看到的，这些子空间可能非常有趣。

到目前为止，我们只是抽象地讨论杂化弦理论。下一步是实际构造这样的理论；代数几何也由此登场。构造方法有好几种，对我们最重要的是由 Calabi--Yau 三维簇\footnote{眼下，我们把 Calabi--Yau 三维簇定义为第一 Betti 数为零、典范类平凡的光滑、紧致、连通三维簇。稍后我们会允许某些奇点。} $V$ 以及 $V$ 上的\emph{复化 Kähler 类} $\omega=B+iJ$ 所确定的\emph{非线性 $\sigma$ 模型}，简称 $\sigma$ 模型。这里 $B,J\in H^2(V,\mathbb R)$，且 $J$ 是 Kähler 类。\index{非线性sigma模型@非线性 $\sigma$ 模型|hyperpage}

从输入数据 $(V,\omega)$ 出发，有一种几何构造可以得到 $N=2$ 的 SCFT，并且明确赋予各种超对称不同的作用。因此，在这一背景下，$(Q,\overline Q)$ 有一个典范选取。于是，我们的 Hilbert 空间上的 $u(1)\times u(1)$ 表示也得到明确选取。可以计算出：对 $p,q\geq0$，$(Q,\overline Q)$ 的特征子空间为
\begin{equation}\label{eq:eigenspaces}
\begin{aligned}
(p,q)\text{ 特征子空间}&\simeq H^q(V,\wedge^pT_V),\\
(-p,q)\text{ 特征子空间}&\simeq H^q(V,\Omega_V^p).
\end{aligned}
\end{equation}
\hyperref[sec:nonlinear-sigma-models]{附录 B.2} 更完整地解释什么是非线性 $\sigma$ 模型；第 \ref{sec:calabi-yau-physics} 节则说明 Calabi--Yau 条件如何从物理学中产生。

杂化弦理论中最重要的场与 $H^1(V,T_V)$ 和 $H^1(V,\Omega_V^1)$ 中的元素相联系；它们分别对应特征值 $(1,1)$ 和 $(-1,1)$。与这些空间中的元素相应的算子，称为\emph{边缘算子}。它们之所以重要，部分原因在于它们与来自 $(V,\omega)$ 的 $\sigma$ 模型的模密切相关。直观地说，同时改变 $V$ 的复结构和复化 Kähler 类 $\omega=B+iJ$，便得到 SCFT 的模；不过，SCFT 模空间还存在一些额外的离散等同，暂且不予考虑。读者或许熟悉复模，但“复化 Kähler 模”这一概念可能比较陌生。我们将在第 \ref{sec:math-preview} 节详细介绍。\index{边缘算子|hyperpage}\index{复化Kähler模@复化 Kähler 模|hyperpage}

由此可知，SCFT 模空间有两个叶状结构：其中一个的叶可描述为“$V$ 不变”，另一个则是“$\omega$ 不变”。这就得到下图：
\begin{figure}[H]
 \centering\refstepcounter{equation}\label{eq:scft-foliations}
 \makebox[\textwidth]{\makebox[0pt][l]{(\theequation)}\hfill
 \raisebox{-.5\height}{\includegraphics[width=0.48536\textwidth]{1-2-scft-moduli.png}}\hfill\mbox{}}
\end{figure}

\SourcePage{27}
尽管有这幅图，我们仍要强调：SCFT 模空间\emph{并非}复结构模空间与 Kähler 结构模空间的乘积，即使在局部也不是。事实上，$\omega$ 的 Kähler 模空间可能依赖于 $V$ 的复结构 \SourceCite{Wilson2}。一般而言，只有先固定 $V$ 的复结构，它才得到良好定义。另一方面，\SourceCite{Wilson2} 表明，对足够一般的 Calabi--Yau 三维簇 $V$，$\omega$ 的 Kähler 模与 $V$ 的复结构无关。\hyperref[chap:moduli]{第~\ref*{chap:moduli} 章}将更详细地讨论这些问题。因此，上图虽有助于概念上的理解，却没有反映 SCFT 模空间的微妙之处。

现在，我们对 Calabi--Yau 三维簇和复化 Kähler 类如何给出有趣的物理理论已有了较清楚的认识，可以解释镜像对称从何而来了。基本出发点是 $(Q,\overline Q)$ 的符号不确定性。前面说过，$(Q,\overline Q)$ 除符号外唯一确定，但由 $(V,\omega)$ 得到的 $\sigma$ 模型却作出了一个十分具体的选取。如果将 $Q$ 改为 $-Q$，而保持 $\overline Q$ 不变，那么 $(p,q)$ 与 $(-p,q)$ 特征子空间就会互换；依照式~\eqref{eq:eigenspaces}，这会交换 $H^q(V,\wedge^pT_V)$ 与 $H^q(V,\Omega_V^p)$。一般来说，这两个向量空间的维数不同，所以对同一个模型不可能如此交换。然而，从物理学的观点看，这种符号改变是合理的。这种不对称性暗示：符号改变或许对应于由另一个不同的配对 $(V^\circ,\omega^\circ)$ 所产生的 $\sigma$ 模型。如果这样的配对存在，就称 $(V,\omega)$ 与 $(V^\circ,\omega^\circ)$ 构成一个\emph{镜像对}。更正式地说，我们有如下来自物理学的定义。\index{镜像对|hyperpage}

\begin{physicsdefinition}\label{def:physics-mirror-pair}
若 $(V,\omega)$ 与 $(V^\circ,\omega^\circ)$ 的 $\sigma$ 模型所诱导的超共形场论同构，而且这些理论的 $N=2$ 超共形表示只相差上述符号改变，就称它们构成一个镜像对。
\end{physicsdefinition}

注意，这不是一个数学定义，因为与 $(V,\omega)$ 相联系的 SCFT 尚未得到严格定义。不过，\hyperref[chap:quantum-cohomology]{第~\ref*{chap:quantum-cohomology} 章}将给出\emph{数学镜像对}的严谨定义，其中会包含镜像对称所预言的许多性质。

下面就来考察这些性质。如果 $(V,\omega)$ 与 $(V^\circ,\omega^\circ)$ 构成镜像对，我们便得到同构的 SCFT。但这在数学上意味着什么？理解镜像对称的数学后果，正是本书的主要目标之一。

要看镜像对称对 $V$ 与 $V^\circ$ 有何启示，先把式~\eqref{eq:eigenspaces} 与特征值的变换 $(p,q)\leftrightarrow(-p,q)$ 结合起来，就得到同构
\begin{equation}\label{eq:mirror-cohomology-tangent}
\begin{aligned}
H^q(V,\wedge^pT_V)&\simeq H^q(V^\circ,\Omega^p_{V^\circ}),\\
H^q(V,\Omega^p_V)&\simeq H^q(V^\circ,\wedge^pT_{V^\circ}).
\end{aligned}
\end{equation}
由于 $V$ 是 Calabi--Yau 簇，它有一个处处非零的全纯三形式 $\Omega$；与 $\Omega$ 作杯积，便给出一个非典范同构 $H^q(V,\wedge^pT_V)\simeq H^q(V,\Omega_V^{3-p})$。$V^\circ$ 也一样。因此，式~\eqref{eq:mirror-cohomology-tangent} 可改写为
\begin{equation}\label{eq:mirror-cohomology-forms}
\begin{aligned}
H^q(V,\Omega_V^{3-p})&\simeq H^q(V^\circ,\Omega_{V^\circ}^p),\\
H^q(V,\Omega_V^p)&\simeq H^q(V^\circ,\Omega_{V^\circ}^{3-p}).
\end{aligned}
\end{equation}
这些同构可以借助 Hodge 菱形得到一个很好的解释。由于 $V$ 是光滑三维簇，它的 Hodge 数 $h^{p,q}(V)=\dim H^q(V,\Omega_V^p)$ 满足对称关系
\[
h^{p,q}(V)=h^{q,p}(V)=h^{3-p,3-q}(V)=h^{3-q,3-p}(V).
\]
又因为 $V$ 是 Calabi--Yau 簇，所以 $b_1(V)=0$ 且 $\Omega_V^3\simeq\mathcal O_V$。这推出 $h^{1,0}(V)=0$、$h^{3,0}(V)=1$。此外，$h^{0,2}(V)=\dim H^2(V,\mathcal O_V)=\dim H^1(V,\mathcal O_V)=$
\SourcePage{28}
$h^{0,1}(V)=0$；其中第二个等号来自 Serre 对偶和 $\Omega_V^3\simeq\mathcal O_V$。因此，$V$ 的 Hodge 菱形为
\[
\begin{array}{ccccccc}
&&&1&&&\\
&&0&&0&&\\
&0&&h^{1,1}(V)&&0&\\
1&&h^{2,1}(V)&&h^{2,1}(V)&&1\\
&0&&h^{1,1}(V)&&0&\\
&&0&&0&&\\
&&&1&&&
\end{array}
\]
现在比较镜像对 $V$ 与 $V^\circ$ 的 Hodge 菱形。式~\eqref{eq:mirror-cohomology-forms} 蕴含 $h^{p,q}(V)=h^{3-p,q}(V^\circ)$，表明 $V^\circ$ 的 Hodge 菱形是 $V$ 的 Hodge 菱形关于一条 $45^\circ$ 直线的反射，即镜像。“镜像对称”这个名称正是由此而来。\index{Hodge菱形@Hodge 菱形|hyperpage}

式~\eqref{eq:mirror-cohomology-tangent} 和~\eqref{eq:mirror-cohomology-forms} 其实只是一系列愈发引人瞩目的镜像对称后果中的头两项。构成镜像对的下一个有趣后果涉及模空间。要理解模空间如何出现，注意式~\eqref{eq:mirror-cohomology-tangent} 给出 $H^1(V,T_V)\simeq H^1(V^\circ,\Omega^1_{V^\circ})$ 和 $H^1(V,\Omega_V^1)\simeq H^1(V^\circ,T_{V^\circ})$。这自然地把 $V$ 的复模空间的切空间与 $\omega^\circ$ 的 Kähler 模空间的切空间等同起来（定义见第 \ref{sec:math-preview} 节）；同样，它把 $\omega$ 的 Kähler 模空间的切空间与 $V^\circ$ 的复模空间的切空间等同起来。因此，$V$ 的复模空间局部同构于 $\omega^\circ$ 的 Kähler 模空间；同样，$V^\circ$ 的复模空间局部同构于 $\omega$ 的 Kähler 模空间。这些局部同构合称为\emph{镜像映射}。\hyperref[chap:moduli]{第~\ref*{chap:moduli} 章}将更详细地研究复模和 Kähler 模，并给出镜像映射的严谨定义。

回想我们那个略有失真的示意图~\eqref{eq:scft-foliations}：$(V,\omega)$ 的 SCFT 模空间有两个叶状结构，其中一个的叶满足“$V$ 不变”，另一个满足“$\omega$ 不变”。因此，若 $(V,\omega)$ 与 $(V^\circ,\omega^\circ)$ 构成镜像对，就得到下列 SCFT 模空间的局部同构：
\begin{figure}[H]
 \centering\refstepcounter{equation}\label{eq:mirror-foliations}
 \makebox[\textwidth]{\makebox[0pt][l]{(\theequation)}\hfill
 \raisebox{-.5\height}{\includegraphics[width=0.86471\textwidth]{1-3-mirror-moduli.png}}\hfill\mbox{}}
\end{figure}
这幅图也说明，谈论 Calabi--Yau 流形 $V$ 的“那个”镜像流形并无意义。给定 $V$ 后，我们可以在 SCFT 模空间的一个叶内任意改变 $\omega$；在镜像一侧，这将使 $\omega^\circ$ 保持不变，而 $V^\circ$ 自由变化。因此，真要说“$V$ 的镜像”，它应当是类 $\omega^\circ$，
\SourcePage{29}
连同 $V^\circ$ 的那些形变所构成的模空间；在这些形变上，$\omega^\circ$ 仍可被理解为一个复化 Kähler 类。

除了上面所讨论的，镜像对的存在还有更多后果，其中并非所有后果都已得到理解。基本思想是：原则上，凡是能够用 SCFT 定义的量，都可以利用 SCFT 的两种不同构造来计算。在最理想的情形下，还可以用 $V$ 与 $V^\circ$ 的几何来计算这些量。注意，$Q$ 的符号发生了改变，而 $H^q(V,\Omega_V^p)$ 又对应于 $H^q(V^\circ,\wedge^pT_{V^\circ})$；因此，预期两个模型中的几何计算会有所不同。最好的例子是式~\eqref{eq:correlation} 中已经提到的 SCFT 关联函数。第 \ref{sec:three-point} 节将更详细地讨论它们；它们也是镜像对称作出枚举预测的关键。我们将看到，在镜像族上进行计算，可以得出关于原来 Calabi--Yau 流形的惊人结论。

在物理学文献中，镜像对称是一个内容丰富的现象。除了前面讨论的非线性 $\sigma$ 模型的镜像对称，一些非几何型的 SCFT 中也观察到了镜像对称，其中包括极小模型和 Landau--Ginzburg 轨形。要解释镜像对称在这些情形下如何起作用，需要对 SCFT 按一个有限群取“轨形”。首先取由不变场组成的子理论；但所得子理论在时间演化下并不稳定，也就是说，它不是幺正的。因此需要增添一些场，得到再次成为 SCFT 的幺正理论。实际的物理过程更微妙，细节见\hyperref[sec:landau-ginzburg-models]{附录 B.4}。如果按经过精心选取的群作用取商，就会得到与出发时相同的物理理论，只是 $Q$ 特征值的符号发生了改变。这种形式的“镜像对称”甚至早于非线性 $\sigma$ 模型中镜像对称的发现。

Calabi--Yau 三维簇存在镜像对称的早期证据，来自加权射影空间中的 Calabi--Yau 超曲面，或它们按有限群取商后所得对象的列表。这些超曲面的 Hodge 数展现出醒目但远非完美的对称性。对于其中一些超曲面，\SourceCite{GPl} 先证明某些 Landau--Ginzburg 理论的镜像对称，再把这些理论与超曲面的 $\sigma$ 模型联系起来，从而证明了相应的镜像对称。正如\hyperref[chap:mirror-constructions]{第 4 章}将说明的，\SourceCite{CdK} 指出，这些加权射影超曲面都属于 Batyrev 的自反多胞形构造 \SourceCite{Batyrev4} 所产生的对象。人们猜想——而且普遍相信——由自反多胞形得到的 Calabi--Yau 三维簇具有镜像对称；更一般地，范围更大的环面完全交类 \SourceCite{Borisov1} 也具有镜像对称。过去几年已有大量证据支持这些看法 \SourceCite{CdGP,Font,Morrison1,CdFKM,HKTY1,CFKM,HKTY2,BK1,AGM1,ES1,BKK,Kontsevich2,MP1,Givental2,Givental4,LLY}。\hyperref[chap:mirror-theorem]{第~\ref*{chap:mirror-theorem} 章}将概述证明镜像定理的两种相关方法；该定理确立 Calabi--Yau 环面完全交与其猜想镜像的某些关联函数相等。\index{自反多胞形|hyperpage}

\section{三点函数}\label{sec:three-point}
式~\eqref{eq:correlation} 定义的关联函数，在 SCFT 中本身就是值得研究的对象。在物理学里，研究粒子的连续世代时，它们会自然出现。最常见的关联函数是\emph{三点函数}，
\SourcePage{30}
它描述来自三个世代——不一定互不相同——的粒子之间的相互作用。在基本粒子物理的标准模型中，粒子的一个\emph{世代}，是指在电磁力、弱核力与强核力作用下具有特定相互作用类型的一组粒子。实验表明，粒子共有三个世代。其中一个世代包含最熟悉的粒子，即电子及与它相伴的中微子，以及构成质子和中子的上夸克与下夸克。已知的另外两个世代则包含更奇特的夸克和轻子。

我们将考察由 Calabi--Yau 三维簇 $V$ 和复化 Kähler 类 $\omega$ 所产生的非线性 $\sigma$ 模型的三点函数。最有趣的三点函数是 Yukawa 耦合，它们来自前一节讨论的边缘算子。边缘算子对应于 $H^1(V,\Omega_V^1)$ 和 $H^1(V,T_V)$，因此需要考虑两类 Yukawa 耦合。\index{Yukawa耦合@Yukawa 耦合|hyperpage}

先看由 $H^1(V,\Omega_V^1)$ 产生的 Yukawa 耦合。对于这个空间中的每个元素，$\sigma$ 模型都对应一个由场组成的 $27$ 维向量空间；这些场构成 $E_6$ 的一个不可约表示。为了把弦理论与现实的物理世界联系起来，人们认为这个向量空间包含一代基本粒子。令 $\omega_i\in H^1(V,\Omega_V^1)$，$i=1,2,3$，分别对应三个世代。那么这些世代之间的 Yukawa 耦合，是从每一代各取一个粒子而形成的、在物理学上重要的三粒子耦合。

在此情形下，用 Feynman 路径积分方法算得 Yukawa 耦合为
\begin{equation}\label{eq:a-yukawa}
\begin{aligned}
\langle\omega_1,\omega_2,\omega_3\rangle
={}&\int_V\omega_1\wedge\omega_2\wedge\omega_3\\
&+\sum_{\substack{\beta\ne0\\\beta\text{ 有效}}}n_\beta
 \left(\int_\beta\omega_1\right)
 \left(\int_\beta\omega_2\right)
 \left(\int_\beta\omega_3\right)
 \frac{e^{2\pi i\int_\beta\omega}}{1-e^{2\pi i\int_\beta\omega}}.
\end{aligned}
\end{equation}
求和遍及非零有效同调类 $\beta\in H_2(V,\mathbb Z)$；粗略地说，$n_\beta$ 是同调类为 $\beta$ 的有理曲线的数目。要严格定义 $n_\beta$，必须理解 Gromov--Witten 不变量；\hyperref[chap:gromov-witten]{第~\ref*{chap:gromov-witten} 章}将对此作详细讨论。理论预言，当 $V$ 发生形变时，$n_\beta$ 不变。因此，这个 Yukawa 耦合依赖于 $\omega$，却不依赖于 $V$。\hyperref[chap:quantum-cohomology]{第~\ref*{chap:quantum-cohomology} 章}将说明，式~\eqref{eq:a-yukawa} 中的耦合与 $V$ 的量子上同调环密切相关。

刚才描述的 Yukawa 耦合有时称为\emph{A 模型关联函数}。这个名称来自\hyperref[sec:nonlinear-sigma-models]{附录 B.2} 所介绍的 A 模型。A 模型是 SCFT 的一个“扭曲”版本：某些场局部看来与 $\sigma$ 模型中的场相同，但从整体来看，它们成为原来所在丛的某种扭曲的截面。这些耦合与 A 模型中相应的三点函数等同。

再简要说明式~\eqref{eq:a-yukawa} 从何而来。注意，其第一项只是杯积：
\begin{equation}\label{eq:a-cup}
\int_V\omega_1\wedge\omega_2\wedge\omega_3.
\end{equation}
从物理学的观点看，由式~\eqref{eq:a-cup} 走向 A 模型关联函数~\eqref{eq:a-yukawa} 时，发生了两件事。首先有一个\emph{非重整化定理}：从微扰的观点看，不需要任何量子修正。正如上一节所说，这是超对称的一个良好后果。不过，还须考虑世界面上的非微扰修正；
\SourcePage{31}
在这里，这些修正来自\emph{全纯瞬子}。全纯瞬子是从紧致 Riemann 曲面 $\Sigma$ 到 $V$ 的非常值全纯映射。第 \hyperref[chap:gromov-witten]{\ref*{chap:gromov-witten}}、\hyperref[chap:quantum-cohomology]{\ref*{chap:quantum-cohomology}} 章更仔细地处理 A 模型关联函数时，我们会看到，$\Sigma$ 可以有结点奇点，也可以有不止一个分支。在 A 模型关联函数中，只须考虑 $\Sigma$ 的亏格为零的瞬子。粗略地说，式~\eqref{eq:a-yukawa} 的 $n_\beta$ 计数的正是它们。按照\hyperref[chap:gromov-witten]{第~\ref*{chap:gromov-witten} 章}的术语，我们称 $n_\beta$ 为\emph{瞬子数}。\index{瞬子数|hyperpage}

第二类需要考虑的 Yukawa 耦合来自 $H^1(V,T_V)$。此时，该空间中的元素对应于 $E_6$ 的共轭 $27$ 维表示。取 $\theta_i\in H^1(V,T_V)$，$i=1,2,3$，分别对应三个世代，便得到相应的 Yukawa 耦合：
\begin{equation}\label{eq:b-yukawa}
\langle\theta_1,\theta_2,\theta_3\rangle
=\int_V\Omega\wedge(\theta_1\cdot\theta_2\cdot\theta_3\cdot\Omega),
\end{equation}
这里 $\Omega$ 是 $V$ 上的一个全纯三形式。表达式 $\theta_1\cdot\theta_2\cdot\theta_3\cdot\Omega$ 由以下复合映射定义：
\[
S^3H^1(V,T_V)\otimes H^0(V,\Omega_V^3)
\longmapsto H^3(V,\wedge^3T_V\otimes\Omega_V^3)
\simeq H^3(V,\mathcal O_V)\simeq H^{0,3}(V).
\]
也可以把式~\eqref{eq:b-yukawa} 理解为
\[
\int_V\Omega\wedge
\bigl(\nabla_{\theta_1}\nabla_{\theta_2}\nabla_{\theta_3}\Omega\bigr),
\]
其中 $\nabla$ 是 Gauss--Manin 联络。\footnote{\SourceCite{CdGP} 中的 Yukawa 耦合就是以这种形式出现的。正如命题~\ref{prop:normalized-yukawa-pf} 将说明的，Hodge 理论促使我们在 Yukawa 耦合的定义中再加一个负号。\SourceCite{Cd} 还讨论了加负号的物理理由。由于 $\Omega$ 可以乘以任意常数，眼下符号并不重要；但在\hyperref[subsec:normalized-yukawa-mirror]{第~\ref*{subsec:normalized-yukawa-mirror} 小节}，我们会固定 $\Omega$ 的一种选取，那时就必须留意符号。我们对归一化 Yukawa 耦合的定义将包含这个负号。}注意，$\langle\theta_1,\theta_2,\theta_3\rangle$ 此时尚未得到良好定义，因为 $\Omega$ 可以乘以任意常数。不过，存在一种自然的归一化方法，后文将予介绍。显然，这个 Yukawa 耦合不依赖于复化 Kähler 类 $\omega$，却依赖于 $V$ 的复结构，因为 $\Omega$ 是全纯三形式。

式~\eqref{eq:b-yukawa} 中的 Yukawa 耦合有时称为\emph{B 模型关联函数}，因为它与另一种扭曲理论——\hyperref[sec:nonlinear-sigma-models]{附录 B.2} 中介绍的 B 模型——里相应的三点函数相同。

要解释式~\eqref{eq:b-yukawa} 的来历，我们仿照 A 模型的情形，从“杯积”
\begin{equation}\label{eq:b-cup}
\int_V\Omega\wedge(\theta_1\cdot\theta_2\cdot\theta_3\cdot\Omega)
\end{equation}
出发。这里同样适用\IndexLocation{indexloc:0944:31}非重整化定理，也有与前面相同的全纯瞬子。关键在于，瞬子通过杯积发挥作用。在 A 模型中，$\omega_i$ 经由 $\int_\beta\omega_i$ 出现在式~\eqref{eq:a-yukawa} 中；这可看成杯积 $g_\beta\smile\omega_i\in H^6(V,\mathbb C)\simeq\mathbb C$，其中 $g_\beta\in H^4(V,\mathbb C)=H^{2,2}(V)$ 是 $\beta$ 的 Poincaré 对偶。而在 B 模型中，$\theta_i\in H^1(V,T_V)\simeq H^1(V,\Omega_V^2)=H^{2,1}(V)$，且 $g_\beta\in H^{2,2}(V)$；因此它们的杯积 $g_\beta\smile\theta_i$ 属于 $H^{4,3}(V)=0$。于是，瞬子不会与式~\eqref{eq:b-cup} 相互作用。这就是式~\eqref{eq:b-cup} 等于 B 模型关联函数~\eqref{eq:b-yukawa} 的一个粗略理由。

\SourcePage{32}
式~\eqref{eq:b-yukawa} 中没有瞬子修正，这一点极为重要。它意味着我们能够用代数几何的标准方法，精确计算 B 模型的 Yukawa 耦合。\hyperref[chap:hodge-yukawa]{第 5 章}将解释具体做法。

现在设 $(V,\omega)$ 与 $(V^\circ,\omega^\circ)$ 构成镜像对，也就是说，与 $(V,\omega)$ 相联系的 $\sigma$ 模型，在适当改变符号后，给出的 SCFT 与 $(V^\circ,\omega^\circ)$ 的 $\sigma$ 模型相同。于是自然得到同构 $H^1(V,\Omega_V^1)\simeq H^1(V^\circ,T_{V^\circ})$。又因为三点函数是 SCFT 的内在对象，应当可以推出：来自 $(V,\omega)$ 上 $\sigma$ 模型的 A 模型关联函数，与来自 $(V^\circ,\omega^\circ)$ 上 $\sigma$ 模型的 B 模型关联函数，在适当归一化后相等。这是镜像对称的一个主要数学后果。由于镜像映射也参与其中，实际的等同细节要更复杂一些；但 A 模型与 B 模型关联函数相等，显然是合理的预期。

关联函数的性质也与镜像对称相吻合。我们已经看到，与 $(V,\omega)$ 相联系的 A 模型关联函数依赖于 $\omega$，却不依赖于 $V$。根据镜像对称以及上面讨论的模空间局部对应，与 $(V^\circ,\omega^\circ)$ 相联系的 B 模型关联函数应当依赖于 $V^\circ$，却不依赖于 $\omega^\circ$。如前所述，B 模型确实具有这一性质。

在五次三维簇的例子中我们将看到，\SourceCite{CdGP} 正是利用 A 模型的三点函数与其镜像的 B 模型三点函数相等，才对五次三维簇上的有理曲线数目作出了一些非凡预测。下一章将详细讨论这个例子。一般而言，镜像对称的一些最重要的后果，来自下面两个事实的结合：
\begin{itemize}
\item A 模型与 B 模型的关联函数相等；
\item B 模型的函数可以精确计算。
\end{itemize}
二者合起来，就能计算 Calabi--Yau 三维簇上的 Gromov--Witten 不变量，并由此得到丰富的枚举信息。当然，这一切的前提是我们能够用数学证明：这些物理后果确实成立。

刚才介绍的途径可以看作研究镜像对称后果的“经典”方法，其中 A 模型与 B 模型的关联函数是主要对象。不过，在镜像对称发现后的那些年里，研究重点略有转移。特别是，人们对 B 模型关联函数与 Gauss--Manin 联络之间的关系理解得更深了；在 A 模型一侧，这推动了量子上同调及 A 型 Hodge 结构变分的发展。Givental 以及 Lian、Liu、Yau 关于镜像定理的工作，又引入了其他值得研究的新对象。我们将看到，等变上同调与局部化在镜像定理的证明中起着重要作用。后面的章节会讨论所有这些思想。

最后，再谈谈式~\eqref{eq:a-yukawa} 与~\eqref{eq:b-yukawa}。二者看似不对称，因为前者远比后者复杂。好在镜像对称很容易解释这种差异。上面说过，A 模型关联函数应随复化 Kähler 类 $\omega$ 变化，而 B 模型关联函数应随 $V$ 的复结构变化。式~\eqref{eq:a-cup} 中的杯积是纯拓扑不变量；所以，若要得到对 $\omega$ 的非平凡依赖，显然还需要其他项。另一方面，
\SourcePage{33}
式~\eqref{eq:b-cup} 中的杯积已经依赖于复结构，因为 $\Omega$ 是 $V$ 上的全纯三形式。因此，不再需要额外的项也是可能的。

\begingroup\linespread{1.22}\selectfont
\section{为什么是 Calabi--Yau 流形？}\label{sec:calabi-yau-physics}

接下来的任务，是解释镜像对称为什么只适用于 Calabi--Yau 流形，尤其是 Calabi--Yau 三维簇。要完整说明非线性 $\sigma$ 模型为什么需要 Calabi--Yau 流形，我们必须大幅转入物理学。\hyperref[sec:nonlinear-sigma-models]{附录 B.2} 和 \hyperref[sec:conformal-field-theories]{B.3} 将对此作部分介绍。这里先概述其中一些想法。出发点是假设时空应当是一个带有号型为 $(9,1)$ 的半 Riemann 度量的十维流形。这个流形局部应为乘积 $M_{3,1}\times V$，其中 $M_{3,1}$ 是狭义相对论中的通常时空，$V$ 是一个紧致的六维 Riemann 流形。基本直观是：$V$ 太小，以致在宏观尺度上看不见，但它对理论的量子方面却不可或缺。若时空维数不是十维，本节前面讨论的共形对称性就无法在量子化过程中保持下来。

如往常一样，其中的物理学并不简单，还涉及一些陌生术语。思路是先考察从物理学观点看更为基本的其他理论，再由此接近非线性 $\sigma$ 模型。具体来说，从低能极限下的一个 $N=1$ 超引力理论出发，先去掉费米场；随后，超对称和其他相容性要求对 Riemann 流形 $V$ 上的度量施加强约束。特别地，可得到该度量的和乐群等于 $SU(3)$。这有一些很好的后果：
\begin{itemize}
\item \textbf{微分几何方面：}$V$ 上的度量是 Ricci 平坦的，即其 Ricci 曲率张量处处为零。这蕴含 $b_1(V)=0$。
\item \textbf{代数几何方面：}$V$ 具有一个复结构，使得 $c_1(V)=0$，而且该度量对于这一复结构是 Kähler 的。
\end{itemize}
由此可见，$V$ 正是我们所需要的 Calabi--Yau 流形。

到目前为止，在时空中我们只有 $N=1$ 超对称。经过进一步处理，这个理论可以重新解释为世界面上的 $N=2$ SCFT；不过，为了保持超共形不变性，还须形变上述 Ricci 平坦度量。我们尚不知道这种形变的精确描述——这仍是物理学中的一个开放问题——但可以证明新度量与旧度量属于同一个 Kähler 类。因此，得到的仍然是一个 Calabi--Yau 三维簇。

如果读者想更完整地了解从物理学观点看 Ricci 曲率如何出现，并查阅相关文献，可参阅 \SourceCite{Huebsch} 的第 0 章。

\section{镜像对称的数学}\label{sec:math-preview}

从数学的观点看，第 \ref{sec:physics} 节给出的镜像对称表述存在严重问题。例如，$N=2$ SCFT 的定义涉及对所有映射 $\Sigma\to V$ 构成的空间积分，而这类积分尚未得到严格定义。因此，镜像对称的数学证明似乎要求在两个数学上尚未得到良好定义的对象之间——即 $(V,\omega)$ 与 $(V^\circ,\omega^\circ)$ 的 $\sigma$ 模型之间——建立同构。甚至第 \ref{sec:physics} 节示意图里的 $N=2$ SCFT 模空间也没有得到良好定义！

\begingroup\emergencystretch=1em
面对这种情况，数学家该怎么办？一种办法是完全避开 SCFT，转而严谨定义 Kähler 模空间、\SourcePage{34}Gromov--Witten 不变量等对象，再证明它们的行为符合镜像对称的预测。另一种办法是接纳物理学及其直观，以更深入地洞察代数几何，特别是镜像对称对 Calabi--Yau 环面完全交的 Gromov--Witten 不变量所作的预测。第三种办法，是表述一个纯数学版本的镜像对称。例如，\SourceCite{Kontsevich3} 提议：复流形 $V$ 的镜像是某个辛流形 $V^\circ$；镜像对称或可表述为导出范畴的等价，将 $V$ 上的凝聚层与由 $V^\circ$ 的 Lagrange 子流形构成的范畴联系起来。还有一种同样迷人、但目前仍带有一定猜测性质的办法：尝试把镜像流形几何地\emph{构造}为特殊 Lagrange 子流形的模空间 \SourceCite{SYZ}。\footnote{一般而言，辛几何在镜像对称中扮演重要角色，不过本书将更多地关注代数几何方面。}实际上，所有这些办法都得到了应用，这正是镜像对称如此令人兴奋的原因。\par\endgroup
本书将主要采用第一种办法，同时不时对相关物理学加以评论。因此，我们的目标，是讨论理解第 \ref{sec:physics}、\ref{sec:three-point} 节的数学方面所涉及的代数几何。与物理理论不同，我们将研究任意维数的 Calabi--Yau 流形，而非仅限于三维。然而，到\hyperref[chap:quantum-cohomology]{第~\ref*{chap:quantum-cohomology} 章}尝试给镜像对一个数学定义时，我们在 Calabi--Yau 三维簇的情形下将取得最好的结果。

先从严格定义 Calabi--Yau 开始。由于\hyperref[chap:mirror-constructions]{第 4 章}将研究的镜像对称构造有时会产生奇异簇，我们需要一个允许某些类型奇点的定义。

\begin{definition}\label{def:calabi-yau-variety}
一个 $d$ 维正规紧致簇 $V$ 称为\emph{Calabi--Yau 簇}，如果它满足：
\begin{enumerate}[(i)]
\item $V$ 至多具有 Gorenstein 典范奇点；
\item $V$ 的对偶化层平凡，即 $\widehat\Omega_V^d\simeq\mathcal O_V$；
\item $H^1(V,\mathcal O_V)=\cdots=H^{d-1}(V,\mathcal O_V)=\{0\}$。
\end{enumerate}
如果 $V$ 还至多具有 Gorenstein、$\mathbb Q$ 析因的终端奇点，则称 $V$ 为\emph{极小 Calabi--Yau 簇}。
\end{definition}
\hyperref[app:singular-varieties]{附录 A} 将回顾对偶化层 $\widehat\Omega_V^d$，以及 Gorenstein 奇点、典范奇点和终端奇点的定义。第 \ref{sec:physics} 节考虑的 Calabi--Yau 三维簇显然符合这里的定义。

第 \ref{sec:physics} 节用到的其他一些术语也与各种模有关。给定流形 $V$ 上所有复结构构成的空间，是代数几何中熟知的对象；Kähler 模的概念则可能不那么熟悉。回想第 \ref{sec:physics} 节：我们有一个 Calabi--Yau 三维簇 $V$，以及一个复化 Kähler 类 $\omega=B+iJ\in H^2(V,\mathbb C)$，其中 $J$ 是 Kähler 类。不过，如果用一个整上同调类改变 $\omega$，物理理论不会改变，因为非线性 $\sigma$ 模型的定义对映射 $\phi:\Sigma\to V$ 只用到 $\exp\bigl(2\pi i\int_\Sigma\phi^*(\omega)\bigr)$。\footnote{从 A 模型关联函数~\eqref{eq:a-yukawa} 中也可以看出这一点。完整细节见\hyperref[sec:nonlinear-sigma-models]{附录 B.2}。}当 $\omega$ 改变一个 $H^2(V,\mathbb Z)$ 中的元素时，这个量保持不变。因此，在定义 Kähler 模时，应当对 $H^2(V,\mathbb Z)$ 的像取商。由此得到下面的定义；和前面一样，它也允许某些奇点。

\SourcePage{35}
\begin{definition}\label{def:kahler-moduli}
设 $V$ 是一个射影轨形，且 $h^{2,0}(V)=0$。那么：
\begin{enumerate}[(i)]
\item \emph{Kähler 锥} $K(V)\subset H^2(V,\mathbb R)=H^{1,1}(V,\mathbb R)$，是所有 Kähler 类构成的子集；
\item \emph{复化 Kähler 空间}是商空间
\[
K_{\mathbb C}(V)=\{\omega\in H^2(V,\mathbb C):\operatorname{Im}(\omega)\in K(V)\}/\operatorname{im}H^2(V,\mathbb Z),
\]
其中 $\operatorname{im}H^2(V,\mathbb Z)$ 是自然映射 $H^2(V,\mathbb Z)\to H^2(V,\mathbb C)$ 的像；
\item \emph{复化 Kähler 模空间}是商空间 $K_{\mathbb C}(V)/\operatorname{Aut}(V)$。
\end{enumerate}
\end{definition}
事实上，在 SCFT 中出现的 Kähler 模空间\footnote{讨论 Kähler 模空间时，只要上下文足以表明含义，我们常省去“复化”这个限定语。}与这里的定义略有不同。我们定义的 Kähler 模空间会受到量子修正，其性质因此会发生细微变化。特别是，理论可能并非对所有 $\omega$ 都收敛，所以我们不得不把注意力限制在 $\omega=B+iJ$ 中 $J$ 充分正的复化 Kähler 类上。尽管如此，这里定义的较大空间在数学上仍然有趣，也是我们的主要研究对象之一。\hyperref[chap:moduli]{第~\ref*{chap:moduli} 章}将讨论其中一些微妙之处。

\hyperref[app:singular-varieties]{附录 A} 将回顾轨形及轨形上 Kähler 类的定义。因为 $h^{2,0}=0$，Kähler 锥 $K(V)$ 是 $H^2(V,\mathbb R)$ 中的开凸锥。这说明复化 Kähler 空间 $K_{\mathbb C}(V)$ 是一个性质良好的对象。但是，若要确定 Kähler 模空间 $K_{\mathbb C}(V)/\operatorname{Aut}(V)$ 的结构，还须知道\IndexLocation{indexloc:0088:35}自同构群 $\operatorname{Aut}(V)$ 如何作用于 Kähler 锥。\hyperref[chap:moduli]{第~\ref*{chap:moduli} 章}将回到这一主题。

注意，定义~\ref{def:calabi-yau-variety} 与定义~\ref{def:kahler-moduli} 似乎涉及略微不同的奇点类型。幸运的是，\SourceCite{Reid1} 的一个结果表明，Gorenstein 轨形具有\IndexLocation{indexloc:1259:35}典范奇点。这意味着，凡满足定义~\ref{def:calabi-yau-variety} 第 (ii)、(iii) 项的轨形，都自动是 Calabi--Yau 簇。因此，对于镜像对称，Calabi--Yau 轨形是一类自然的研究对象。

\section{后续内容}\label{sec:whats-next}

下面介绍接下来的十章将如何把我们从这里引向镜像定理的证明。\hyperref[chap:quintic]{第 2 章}首先仔细描述光滑五次三维簇 $V\subset\mathbb P^4$ 的镜像。沿着第 \ref{sec:three-point} 节概述的策略，我们将对 $V$ 上给定次数的有理曲线数目作出一些明确预测。在后续章节中，随着数学背景逐步建立，我们还会多次回到这个例子，并最终论证\hyperref[chap:quintic]{第 2 章}中出现的全部计算。

人们很快意识到，要推广五次三维簇的例子，环面几何有重要作用。因此，\hyperref[chap:toric]{第 3 章}将探讨描述环面簇的各种方式，\hyperref[chap:mirror-constructions]{第 4 章}再介绍已知的镜像对称构造，其中许多构造都用到环面几何。这会为我们提供大量预期具有镜像对称的例子，成为检验数学理论的良好场所。还应提到，有些物理理论也明确使用环面簇，例如\hyperref[sec:gauged-linear-sigma-models]{附录 B.5} 将介绍的 \SourceCite{Witten5} 中的\IndexLocation{indexloc:0354:35}规范线性 $\sigma$ 模型。\SourceCite{MP1} 提供了在这种模型背景下支持镜像对称的证据。

下一个任务是描述并计算 B 模型的 Yukawa 耦合，这将在\hyperref[chap:hodge-yukawa]{第 5 章}完成。事实表明，为了找到计算这一耦合所需的正确坐标，
\SourcePage{36}
计算这个 Yukawa 耦合，需要充分理解模空间在某些边界点附近的情形。因此，\hyperref[chap:moduli]{第~\ref*{chap:moduli} 章}将研究这些模空间的结构及其紧化。正如前面提到的，我们也会考察 Kähler 模。事实上，Kähler 模的若干基本事实，会帮助我们理解通常的模空间的紧化应当是什么样子。\hyperref[chap:gromov-witten]{第~\ref*{chap:gromov-witten} 章}将讨论 Gromov--Witten 不变量。辛几何范畴和代数范畴中都已有相应定义；任一种定义都可用于从数学上定义 A 模型的 Yukawa 耦合。

到\hyperref[chap:gromov-witten]{第~\ref*{chap:gromov-witten} 章}末，我们将具备精确表述“经典”镜像对称所需的一切：它断言某些关联函数在镜像映射下彼此相容。不过，从\hyperref[chap:quantum-cohomology]{第~\ref*{chap:quantum-cohomology} 章}起，我们将探索对这个主题更深层的理解。新引入的两个要素是量子上同调和 A 型 Hodge 结构变分。前者可视为对 Gromov--Witten 不变量的代数与枚举后果的系统展开；后者则自然地产生于量子上同调。基本想法是：镜像对称真正涉及两个 Hodge 结构变分在镜像映射下的同构——一个来自 B 模型，即由复模产生的通常 Hodge 结构变分；另一个来自 A 模型，即 A 型 Hodge 结构变分。\hyperref[chap:quantum-cohomology]{第~\ref*{chap:quantum-cohomology} 章}末将以这种形式表述镜像定理。所要求的关联函数相等，将直接从 Hodge 结构变分的这一同构推出。

用 Hodge 结构变分表述的镜像定理仍是一个开放问题；不过，Givental \SourceCite{Givental2,Givental4} 以及 Lian、Liu、Yau \SourceCite{LLY} 的近期工作，已朝证明这种形式的定理迈出实质性的一步。\hyperref[chap:mirror-theorem]{第~\ref*{chap:mirror-theorem} 章}将讨论这些论文的思想。我们将看到，这需要引入新的技术和新的研究对象，尤其是等变上同调与局部化，它们在证明中扮演重要角色。\hyperref[chap:localization]{第~\ref*{chap:localization} 章}回顾若干基本定义和定理，并运用这些方法证明一些关于 Gromov--Witten 不变量的有趣结果。随后，\hyperref[chap:quantum-differential-equations]{第~\ref*{chap:quantum-differential-equations} 章}讨论 Gromov--Witten 不变量的一种推广，即引力关联数。这些不变量使我们能够描述 A 模型联络的平坦截面。我们还将定义 Givental 的 $J$ 函数，并解释它与量子微分方程的关系。最后，\hyperref[chap:mirror-theorem]{第~\ref*{chap:mirror-theorem} 章}讨论 \SourceCite{Givental2,Givental4} 与 \SourceCite{LLY} 所陈述的镜像定理。\hyperref[chap:quintic]{第 2 章}结尾会简短预览其中涉及的内容，\hyperref[chap:mirror-theorem]{第~\ref*{chap:mirror-theorem} 章}再给出五次三维簇情形的完整细节。

这些章节涉及大量数学内容，从头到尾逐章读下去并非轻松之事。我们建议读者先从\hyperref[chap:quintic]{第 2 章}开始。此后，可依个人兴趣选择几条不同的路线。前言给出了阅读指引，浏览目录也会有所帮助。我们希望，一旦读者开始阅读，这个主题本身的吸引力便会引领读者继续前行。镜像对称的故事引人入胜，而看清各个部分如何拼合起来，也是一件有趣的事。

\par\endgroup
